\subsection{规则函数}

定义 2. 我们称从区间 \(I \subset R\) 到集 E 中的映射 f 为阶梯函数，如果存在 I 分成有限个区间 \(J_{k}\) 的一个分割，使 f 在每个 \(J_{k}\) 上取常值。

设 \((a_{i})_{0\leqslant i\leqslant n}\) 是由诸 \(\mathrm{J}_k\) 的互异端点组成的严格递增序列；由于诸 \(\mathrm{J}_k\) 两两不相交，其中每一个或者退化为一个单点 \(a_{i}\)，或者是以两个相邻的点 \(a_{i}\)、\(a_{i + 1}\) 为端点的区间；此外，由于 I 是诸 \(\mathrm{J}_k\) 之并，点 \(a_0\) 是 I 的左端点，\(a_{n}\) 是 I 的右端点。因此 I 上的每个阶梯函数都可以刻画为在诸开区间 \(]a_{i}\)、\(a_{i+1}[ (0 \leqslant i \leqslant n-1)\) 的每一个上取常值的函数，其中 \((a_{i})_{0 \leqslant i \leqslant n}\) 是 I 的点的严格递增序列，\(a_{0}\) 是 I 的左端点，\(a_{n}\) 是 I 的右端点。

命题 2. 定义在 I 上、取值于 R 上的向量空间 E 中的阶梯函数所成之集，是 I 到 E 中的一切映射所成的向量空间 \(\mathcal{F}(\mathrm{I};\mathrm{E})\) 的向量子空间 E。

事实上，设 \(f\) 与 \(g\) 是两个阶梯函数，\((\mathrm{A}_i)\) 与 \((\mathrm{B}_j)\) 是 I 分成有限个区间的两个分割，使得 \(f\)（相应地，\(g\)）在诸 \(\mathrm{A}_i\)（相应地，\(\mathrm{B}_j\)）的每一个上取常值；无论实数 \(\lambda, \mu\) 如何，显然 \(\lambda f + \mu g\) 在诸非空区间 \(\mathrm{A}_i \cap \mathrm{B}_j\) 的每一个上取常值，并且这些区间构成 I 的一个分割。

推论。向量子空间 E 由区间的特征函数生成。

现在考虑 E 是 R 上赋范空间的情形；这时立即可见，以 a, b（a \textless{} b）为端点的区间 J 的特征函数具有原函数，即当 \(x \leqslant a\) 时等于 a、当 \(a \leqslant x \leqslant b\) 时等于 x、当 \(x \geqslant b\) 时等于 b 的那个函数。于是命题 2 的推论表明，每个取值于 E 中的阶梯函数都具有原函数。

我们现在可以应用 \(n^{\circ}\) 2 中所述的方法。

定义 3. 我们称定义在区间 I 上、取值于 R 上的完备赋范空间 E 中的向量函数为规则函数，如果它是阶梯函数在 I 的每个紧子集上的一致极限。

换言之，规则函数是阶梯函数子空间 E 在 \(\mathcal{F}_{c}(\mathrm{I};\mathrm{E})\) 中的闭包的元素；\(\overline{E}\) 是 \(\mathcal{F}_{c}(\mathrm{I};\mathrm{E})\) 的向量子空间，并且由于 \(\mathcal{F}_{c}(\mathrm{I};\mathrm{E})\) 完备，\(\overline{E}\) 也完备；换言之，若一个函数是规则函数在 I 的每个紧子集上的一致极限，则它在 I 上是规则的。为使 f 在区间 I 上是规则的，必须且只需它在含于 I 的每个紧区间上的限制是规则的。

II, p.~53 的推论 1 表明：

定理 2. 区间 I 上的每个规则函数都在 I 上具有原函数。

我们把 II, p.~4 的定义 3 变换为另一个等价的定义：

定理 3. 为使定义在区间 I 上、取值于 R 上的完备赋范空间 E 中的向量函数 f 是规则的，必须且只需它在 I 的每个内点处具有右极限与左极限，并且当 I 的左端点与右端点属于 I 时，它在左端点处具有右极限、在右端点处具有左极限。因而 f 在 I 中的不连续点所成之集是可数的。

由于每个区间都是可数个紧区间的并，可以只就 I 是紧区间的情形证明定理 3，设 I = {[}a, b{]}。

\(1^{\circ}\) 此条件是必要的。设 \(\mathbf{f}\) 是规则的，并设 \(x\) 是 I 的异于 \(b\) 的一点。依假设，对每个 \(\varepsilon > 0\) 存在阶梯函数 \(\mathbf{g}\)，使得 \(\| \mathbf{f}(z) - \mathbf{g}(z) \| \leqslant \varepsilon\) 对每个 \(z \in \mathrm{I}\) 成立；由于 \(\mathbf{g}\) 在点 \(x\) 处具有右极限，存在 \(y\)，使得 \(x < y \leqslant b\)，并且对每对点 \(z\)、\(z'\)（在区间 \([x, y]\) 中）有 \(\| \mathbf{g}(z) - \mathbf{g}(z') \| \leqslant \varepsilon\)，因而有 \(\| \mathbf{f}(z) - \mathbf{f}(z') \| \leqslant 3\varepsilon\)；这就证明（依柯西准则）\(\mathbf{f}\) 在点 \(x\) 处具有右极限。同理可证 \(\mathbf{f}\) 在 I 的每个异于 \(a\) 的点处具有左极限。

\(2^{\circ}\) 此条件是充分的。设它成立；对每个 \(x \in I\) 存在开区间 \(\mathrm{V}_x = ]c_x\)、\(d_x[\)，包含 \(x\)，并且使在 \(I\) 与诸开区间 \(]c_x\)、\(x[i, ]x\)、\(d_x[\) 中每一个的交（当此交非空时）上 \(\mathbf{f}\) 的振幅为 \(\leqslant \varepsilon\)。由于 \(I\) 紧，存在有限个点 \(x_i\)（在 \(I\) 中），使诸 \(\mathrm{V}_{x_i}\) 构成 \(I\) 的一个覆盖；设 \((a_k)_{0 \leqslant k \leqslant n}\) 是把由点 \(a\)、\(b\) 以及 \(x_i\)、\(c_{x_i}\) 与 \(d_{x_i}\)（属于 \(I\) 者）所组成的有限集的点按递增次序排列所得的序列；诸区间 \(]a_k\)、\(a_{k+1}[\)（\(0 \leqslant k \leqslant n-1\)）中的每一个都含于某个区间 \(]c_{x_i}\)、\(x_i[\) 或 \(]x_i\)、\(d_{x_i}[\)，故 \(\mathbf{f}\) 在其上的振幅为 \(\leqslant \varepsilon\)；设 \(\mathbf{c}_k\) 是 \(\mathbf{f}\) 在 \(]a_k\)、\(a_{k+1}[\) 上的一个值；令 \(\mathbf{g}(a_k) = \mathbf{f}(a_k)\)（对 \(0 \leqslant k \leqslant n\)），并令 \(\mathbf{g}(x) = \mathbf{c}_k\) 对一切 \(x \in ]a_k\)、\(a_{k+1}[\)（\(0 \leqslant k \leqslant n-1\)），便定义了一个阶梯函数 \(\mathbf{g}\)，使得 \(\| \mathbf{f}(z) - \mathbf{g}(z) \| \leqslant \varepsilon\) 在 \(I\) 上成立；所以 \(\mathbf{f}\) 在 \(I\) 上是规则的。

最后，我们来证明若 \(\mathbf{f}\) 在 I 上是规则的，则它的不连续点所成之集是可数的。对每个 \(n > 0\) 存在阶梯函数 \(\mathbf{g}_n\)，使得在 I 上 \(\| \mathbf{f}(x) - \mathbf{g}_n(x)\| \leqslant 1 / n\)；由于序列 \((\mathbf{g}_n)\) 在 I 上一致收敛于 \(\mathbf{f}\)，我们看出 \(\mathbf{f}\) 在诸 \(\mathbf{g}_n\) 都连续的每个点处连续（Gen.~Top., X, p.~281, 推论 1）；但由于 \(\mathbf{g}_n\) 除去有限集 \(\mathrm{H}_n\) 的点外连续，故 \(\mathbf{f}\) 在集 \(\mathrm{H} = \bigcup_{n}\mathrm{H}_{n}\) 的余集的点处连续，而该余集是可数的。

推论 1. 设 \(\mathbf{f}\) 是 I 上的规则函数；在 I 的每个点处，除去 I 的右端点（相应地，左端点），\(\mathbf{f}\) 的每个原函数都具有等于 \(\mathbf{f}(x+)\) 的右导数（相应地，等于 \(\mathbf{f}(x-)\) 的左导数）；特别地，在每个点 \(x\)（\(\mathbf{f}\) 在其处连续）处，\(\mathbf{f}(x)\) 是其某个原函数的导数。

这是 II, p.~54 的定理 3 与 I, p.~18 的命题 6 的直接推论。

推论 2. 设 \(\mathbf{f}_i\)（\(1 \leqslant i \leqslant n\)）是区间 I 上的 \(n\) 个规则函数，每个 \(\mathbf{f}_i\) 取值于一个完备赋范空间 \(\mathrm{E}_i\) 中（在 \(\mathbf{R}\)（\(1 \leqslant i \leqslant n\)）上）。若 \(\mathbf{g}\) 是子空间 \(\prod_{i=1}^{n} \overline{\mathbf{f}_i(\mathrm{I})}\)（\(\prod_{i=1}^{n} \mathrm{E}_i\) 的）到 \(\mathbf{R}\) 上的完备赋范空间 F 中的连续映射，则复合函数 \(x \mapsto \mathbf{g}(\mathbf{f}_1(x), \mathbf{f}_2(x), \ldots, \mathbf{f}_n(x))\) 在 I 上是规则的。

事实上，它显然满足 II, p.~54 的定理 3 的条件。

于是可以看出，若 \(\mathbf{f}\) 是 I 上的规则向量函数，则实函数 \(x\mapsto \| \mathbf{f}(x)\|\) 也是规则的。此外，I 上的实规则函数构成一个环；

而且，若 \(f\) 与 \(g\) 是两个实规则函数，则 \(\sup(f, g)\) 与 \(\inf(f, g)\) 都是规则的。

注记。1) 若 f 是 I 上的实规则函数，而 g 是包含 \(f(\mathrm{I})\) 的区间上的规则向量函数，则复合函数 \(g \circ f\) 未必是规则的（参见~II, p.~79, 习题 4）。

II, p.~54 的定理 3 的两个特殊情形特别重要：

命题 3. 区间 \(I \subset R\) 上取值于 R 上的完备赋范空间 E 中的每个连续向量函数都是规则的，并且在 I 上具有一个原函数，且它在每一点处都是该原函数的导数。

注记。2) 为证明连续函数具有原函数，可以利用每个（系数在 E 中的）实变量多项式函数都具有原函数这一事实；由于由魏尔斯特拉斯定理（Gen.~Top., X, p.~313, 命题 3），每个连续函数都是多项式在每个紧区间上的一致极限，故 II, p.~52 的定理 1 表明每个连续函数都具有原函数。

\begin{enumerate}
\def\labelenumi{\arabic{enumi})}
\setcounter{enumi}{2}
\tightlist
\item
  前一注记的原理可以不作实质性修改地推广到取值于 \(\mathbf{C}\) 上的完备赋范空间中的复变量向量函数。若 \(\mathrm{U}\) 是 \(\mathbf{C}\) 中的开集，同胚于 \(\mathbf{C}\)，则这样一个向量函数 \(\mathbf{f}\)（定义在 \(\mathrm{U}\) 上）的原函数按定义是 \(\mathrm{U}\) 上的连续函数，具有等于 \(\mathbf{f}\) 的导数（在 \(\mathrm{U}\) 的每一点处）。有了这个定义，II, p.~52 的定理 1 可以不作修改地推广（利用 \(\mathrm{U}\) 的连通性证明 \((\mathbf{g}_{\alpha})\) 关于 \(\mathcal{F}\) 在 \(\mathrm{U}\) 的每个点的一个邻域上一致收敛，由此推出 \((\mathbf{g}_{\alpha})\) 关于 \(\mathcal{F}\) 在 \(\mathrm{U}\) 的每个紧子集上一致收敛；再利用 I, p.~18 的命题 4 完成证明）。因此，每个是多项式在 \(\mathrm{U}\) 的每个紧子集上的一致极限的函数都在 \(\mathrm{U}\) 上具有原函数；这些函数正是 \(\mathrm{U}\) 上称为全纯的函数，我们将在后一册书中详细研究它们。
\end{enumerate}

命题 4. 区间 \(I \subset R\) 上的每个单调实函数 f 都是规则的，并且 f 的每个原函数都在 I 上凸。

事实上，f 满足 Gen.~Top., IV, p.~350 的命题 4 的定理 3 的判别法；命题的第二部分由 II, p.~55 的推论 1 以及 I, p.~27 的命题 5 得出。

注记。4) 不要认为区间 I 上的规则函数是仅在 I 上具有原函数的函数（参见~II, p.~80, 习题 7 与 8）。

